\lecture{2}
\title{C in xv6}

\begin{document}

\subsection{Pre-Lecture Reading}

In C, a number written with a leading zero
is assumed to be in octal.  For example,
\verb|0177| equals \verb|127| in decimal.
\begin{itemize}
    \item More generally,
    a prefix of \verb|0x| indicates hexadecimal,
    and a prefix of \verb|0b| indicates binary.
    \item To zero out the last six bits of \verb|x|,
    one might write \verb|x = x & ~077|,
    where \verb|~| is the one's complement.
    \item To extract \verb|n| bits from \verb|x|,
    the highest at position \verb|p|, one might write the following function:
\begin{verbatim}[c]
unsigned getbits(unsigned x, int p, int n)
{
    return (x >> (p+1-n)) & ~(~0 << n);
}
\end{verbatim}
    Note, in particular, that \verb|~(~0 << n)|
    produces an \verb|n|-bit string with only \verb|1|s.
\end{itemize}
We now provide examples of nontrivial idiomatic conventions in C.
\begin{verbatim}[c]
*p++ = val;     /* push val onto stack */
val = *--p;     /* pop top of stack into val */
\end{verbatim}
Here, \verb|p| always points to the empty space in a stack.
Notice how \verb|++| and \verb|--| act as both
\emph{postfix} and \emph{prefix} operators.

\begin{verbatim}[c]
void strcpy(char *s, char *t)
{
    while (*s++ = *t++)
        ;
}
\end{verbatim}
Here, note that \verb|*s = *t| returns \verb|*t|,
so the loop ends when \verb|*t == \0| because \verb|\0| is false-y.

What makes C good as a programming language
is that there is no work done in ``runtime'';
whereas code written in Python and JavaScript requires
nontrivial effort to compile into machine code,
code written in C maps almost directly into assembly.

\subsection{Process Memory and Compilation}

When a process is run, the process memory looks like
(from \verb|0x000...| to \verb|0xfff...|).
\begin{center}
    \frame{\includegraphics[width=0.5\textwidth]{lecture-02/process-memory}}
\end{center}
Recall that memory in the stack is managed by the stack pointer
(e.g., \verb|ld a0, 8(sp)|),
whereas memory in the heap is managed by explicit calls
(e.g., \verb|malloc|).

Consider, in xv6, the following command:
\begin{verbatim}
make user/_cat
\end{verbatim}
The effect of this command is
to compile the source code in the C file \verb|user/cat.c|
into a binary (executable) file \verb|user/_cat|.

There is a compiler \verb|cc| that can compile C files like
\verb|user/cat.c| into \textbf{object files} like \verb|user/cat.o|.
Object files are plain binary---so \verb|cat user/cat.o| will output
unreadable non-ASCII garbage.  Any object file contains the following:
\begin{itemize}
    \item \verb|.text| Assembly instructions written in binary.
    \item \verb|.data| Initialized global variables declared
    outside of the scope of \verb|main|.
    \item \verb|.bss| Uninitialized global variables
    (not stored as data, but rather an indication of size needed).
    \item \dots a bunch of other stuff \dots
\end{itemize}
Consider the following C file, for example.
\begin{verbatim}[c]
#include "user/user.h"        // (1)

int mode = 5;                 // (2)
char buf[64];                 // (3)

int main(void) {              // (4)
  printf("mode=%d\n", mode);  // (5)
  exit(0);                    // (6)
}
\end{verbatim}
Then lines \verb|(4-6)| contribute to \verb|.text|,
line \verb|(2)| contributes to \verb|.data|,
and line \verb|(3)| contributes to \verb|.bss|.
(Meanwhile, line \verb|(1)| is only relevant to
the compiler \verb|cc| and has no impact on
the content of the object file itself.)

It isn't just \verb|user/cat.c| that needs to be compiled;
the code in \verb|user/cat.c| also relies on these four object files:
\begin{verbatim}
ULIB = $U/ulib.o $U/usys.o $U/printf.o $U/umalloc.o
\end{verbatim}
The \verb|make| command knows these files are relevant
because \verb|Makefile| contains the above line, verbatim.
(All four object files are also the results of
compiling their respective source files via \verb|cc|.)

Finally, a \textbf{linker} \verb|ld| takes all five object files
and merges them into a single binary (executable) file \verb|user/_cat|.

There's just a little more work to do.
When typing \verb|cat| into the xv6 shell,
only the binary present in \verb|fs.img| gets executed---none
of the files in \verb|user/| get acknowledged at all!
Here's how we get the binary \verb|user/_cat| into \verb|fs.img|.
\begin{verbatim}
make user/_cat  # create (or update) files user/cat.o and user/_cat
make fs.img     # add the binary user/_cat to fs.img
make qemu       # boot up the xv6 kernel
$ cat           # runs cat
...
\end{verbatim}
Note that the commands \verb|make user/_cat| and \verb|make fs.img|
are technically optional;
when running \verb|make qemu|, the binary contents of \verb|fs.img|
are automatically updated accordingly.

\begin{remark}
    Object files (e.g., \verb|user/cat.o|)
    and binary files (e.g., \verb|user/_cat|)
    both fall under the broad category of \textbf{ELF files},
    which includes object files, executables,
    shared libraries, and core dumps.
\end{remark}

\subsection{C Examples: Pointers, Structs, Arrays, Strings}

Here's a simple example to demonstrate pointers in use.

\begin{verbatim}[c]
int i = 0x12345678;     // stored as four bytes, in the order 78, 56, 34, 12 (little-endian)
int *p = &i;
char *cp = (char *)p;   // `char` is equivalent to `byte` in C
printf("%x", *cp);      // prints 78, not 12
\end{verbatim}

There are two distinctions between pointers of type
\verb|char*| and \verb|int*| to be aware of:
\begin{itemize}
    \item Because \verb|cp| is of type \verb|char*|,
    dereferencing \verb|cp| returns only one byte (a \verb|char|),
    rather than all four bytes of an \verb|int|.
    \item Incrementing \verb|cp++| would increase \verb|cp| by 1 byte,
    whereas incrementing \verb|p++| would increase \verb|p| by 4 bytes.
    \begin{itemize}
        \item The stack can't see type; it sees all pointers as
        plain 8-byte blocks in memory.
        \item It is the \emph{compiler} that can read types
        and realize (in its assembly compilation) that
        \verb|cp| and \verb|p| increment differently.
    \end{itemize}
\end{itemize}

Here's a more complicated example of pointers in use.

\begin{verbatim}[c]
int g = 3;

int
main(int ac, char **av)
{
    int l = 5;  // local variables don't have a default value

    // reading the address of global and local variables
    int *p, *q;
    p = &g;
    q = &l;
    printf("p %p q %p\n", p, q); // p 0x...01000 q 0x...03FAC

    // reassigning variables using pointers
    *p = 11;
    *q = 13;
    printf("g %d l %d\n", g, l); // g 11 l 13

    // structs
    struct two {
        int a;
        int b;
    } s;
    s.a = 10;
    struct two *ptr = &s;
    printf("%d %d\n", s.a, ptr->a); // 10 10

    // make a pointer to a pointer
    int **pp;
    pp = &p;
    printf("pp %p %p %d\n", pp, *pp, **pp); // pp 0x...03FA0 0x...01000 11

    // make a pointer to code
    int (*f)(int, char **);
    f = &main;
    printf("main: %p\n", f); // main: 0x...092

    // make a pointer to a function in a different file (user/ulib.c)
    int (*g)(int, char **);
    extern void start(void);
    g = (void*) &start;
    printf("main: %p\n", g); // 0x...09A
    exit(0);
}
\end{verbatim}

Notice how the pointer \verb|0x...01000| to \verb|g|
is much smaller than the pointer \verb|0x...03FAC| to \verb|l|.
This is because \verb|g| is a global variable
and lives in the data section of process memory,
whereas \verb|l| is a local variable and lives in the stack.

Also, notice how data is added to the stack by moving \emph{downward};
the variable \verb|p| has a \emph{lower} address than \verb|l|,
because \verb|p| was declared \emph{after} \verb|l| was.

Here's how pointers are used to work with arrays in C.

\begin{verbatim}[c]
int a[3] = {1, 2, 3};
char b[3] = {'a', 'b', 'c'};

int
main(int ac, char **av)
{
    printf("%d\n", a[0]); // 1
    
    a[1] += 1;
    *(a + 2) = 5;
    printf("%d %d\n", a[1], a[2]); // 3 5

    printf("a %p b %p", a, b); // a 0x...01010 b 0x...01000
    exit(0);
}
\end{verbatim}

Here's an example of strings in C.
Recall that strings are arrays of characters that end in
the \textbf{zero byte} given by \verb|\0|.

\begin{verbatim}[c]
char *s = "123";

int
main(int ac, char **av)
{
    char s1[4] = {'1', '2', '3', '\0'};
    printf("s %s s1 %s\n", s, s1); // s 123 s1 123

    // accessing arrays
    printf("%c %c\n", s[0], *s); // 1 1
    printf("%c %c\n", s[2], *(s + 2)); // 3 3

    // reading beyond string end (DON'T DO THIS)
    printf("%d %p %p\n", s1[4], s1, &s1[4]); // 0 0x...03FA8 0x...03FAC
    s1[4] = 'D'; // might be benign, might not
    exit(0);
}
\end{verbatim}

Even though data is added to the stack moving \emph{downward},
notice how \verb|s1| has a \emph{lower} address than \verb|s1[4]|,
even though \verb|s1| comes \emph{before} \verb|s1[4]| in the array.
The elements in a single array are stored in ascending address order,
but arrays declared one after another are stored in descending address order.

(This is required, in a sense, in order for \verb|&s1[4] == s1 + 4| to be true.)

Finally, here's the C code given in the pre-lecture homework,
along with the output when \verb|h| is \verb|0x...01000|.

\begin{verbatim}[c]
struct f {
  int a;
  char b[32];
};

struct g {
  char *c;
  int d[4];
};

struct h {
  struct f *f;
  struct g g;
};

struct h *h;

printf("%p\n", &h->g.d[2]);     // 0x...01018
printf("%p\n", &h->f->a);       // 0x...00000
printf("%p\n", &h->g.c);        // 0x...01008
printf("%p\n", &h->g.c[10]);    // 0x...0000A
\end{verbatim}
The first and third are determined by \verb|h| alone;
the second and fourth are not.

\end{document}
